\section{Threats to Validity}
\label{sec:validity}

Every result in Section~\ref{sec:results} depends on choices that a reader
cannot see in the numbers. This section states them, ordered by how far each
could move a conclusion, and it names the one that could move the most first.

\subsection{The Control-Traffic Subsidy}
\label{sec:subsidy}

Centralized protocols are charged nothing for the uplink that informs the sink
of node state. This is the largest single confound in the study and it favors
five of the ten configurations.

The justification is quantitative. A separate 200-bit status report per node per
round, over a sink distance above 100~m and therefore on the fourth-power
branch, costs roughly $3.6\times10^{-5}$~J per node per round, accumulating to
about a quarter of the entire 1~J budget across the round cap. Charging it would
sink all five centralized protocols for a reason unconnected to whether their
clustering decisions are any good, and the piggyback assumption is standard in
the centralized-clustering literature. But a defensible modeling choice is still
a modeling choice, and its magnitude belongs in the results rather than in a
footnote.

Measured over the full sweep, LEACH spends 12.73\% of its consumed energy on
control traffic against 2.33\% for the fuzzy system, 2.46\% for NSGA-II, 2.38\%
for the map, 3.95\% for the DQN and 3.85\% for the GCN. The subsidy is therefore
worth roughly nine to ten percentage points of the total budget. In packet terms
the gap is starker still, 171{,}644 control packets for LEACH and 573{,}870 for
TEEN against 9{,}657 to 12{,}216 for the centralized five, but the packet gap
overstates the effect, since control packets are a fifth the size of data packets
and most are short-range.

Two consequences follow. First, a claim of the form ``NSGA-II extends first
death by 80\% over LEACH'' conflates clustering quality with an accounting
convention and should not be made without qualification. Second, the reliable
comparisons in this paper are the within-class ones: centralized against
centralized, distributed against distributed, with the chain protocol comparing
cleanly to neither. Section~\ref{sec:results}'s strongest results, the
channel-conditional retraction and the DQN-versus-GCN rotation contrast, are both
within-class and are unaffected.

The clean resolution is an ablation that charges the uplink honestly and re-runs
the sweep. It would convert this caveat into a measurement, and it is the single
most valuable follow-up this codebase supports. It has not been run.

\subsection{Transmit Power Is Bracketed, Not Swept}

Section~\ref{sec:retraction} treats the channel comparison as a result, which it
is, but the design bounds what it establishes. Two configurations bracket a
continuum; they do not trace it. The lossy setting places the sink hop near the
packet-error waterfall and the ideal setting removes loss entirely, so a
conclusion stable across both is stable across that interval. A transmit power
between them, or above the lossy setting, would produce intermediate behavior
this study does not measure. A three-point or four-point sweep over transmit
power would trace the curve properly and would show where the fuzzy system's and
the DQN's first-death advantage over LEACH actually crosses the significance
threshold. That sweep has not been run.

The scale sweep of Section~\ref{sec:scale} partially compensates, because moving
the sink changes how much of the node population sits in the waterfall and
reaches the same boundary from a different direction. It is not a substitute:
field size changes area and channel severity together, so it identifies the
boundary without isolating the parameter.

Related parameters are equally unspecified in the source material and equally
consequential: the path-loss exponent, the shadowing standard deviation, the
noise floor and the retry budget. We used conventional values for 2.4~GHz
outdoor short-range links~\cite{rappaport2002wireless,ieee802154} and published
the resulting packet-error curve so that the choice is auditable, but we did not
vary them.

\subsection{The Pre-Training Asymmetry}

The DQN and the GCN arrive pre-trained; the other seven protocols do not. Three
distinct caveats follow and all belong in any citation of these results.

Their reported setup time and operation counts are inference costs only.
Training is real, is paid once, and is reported separately. It must not be read
as free.

The seven non-learned protocols receive no equivalent offline preparation. This
is a fair comparison of a deployed system and it is not an apples-to-apples
algorithmic comparison. Both statements are true simultaneously, and which one
matters depends on the question being asked.

Generalization is now measured rather than assumed, but only in one direction.
The evaluation seeds are disjoint from the training seeds, so the headline
result is not memorization. The scale sweep additionally carries the frozen
$N = 100$, $100\times100$~m weights into eight cells the policies never saw, and
neither policy collapses; the GCN's deficit against LEACH even narrows as the
field grows. What that does not establish is how these policies would compare
against retrained versions of themselves, or how they behave under a sink
placement that is not geometrically similar to the training one. Every cell in
the sweep keeps the base station at $(W/2, 1.5W)$.

\subsection{Modeling Simplifications}

The radio model accounts for transmit, receive and fusion energy and nothing
else. There is no idle-listening cost, no sleep-wake transition energy, no radio
startup and no processor energy. Idle listening frequently dominates real
deployments, and including it would compress the differences between all ten
configurations toward zero. Energy is a linear reservoir with no rate-capacity
effect, no recovery effect and no voltage cutoff.

Fusion is ideal. A head that receives any number of member packets emits exactly
one, so compression is lossless and free in bits. This systematically favors
high-fan-in structures, and PEGASIS most of all: its aggregation ratio of 80
readings per delivered packet is exactly the quantity this assumption idealizes.
The engine enforces the rule symmetrically, so no protocol can exploit it beyond
what its own topology allows, but the assumption still shapes the
chain-versus-cluster comparison more than any other single simplification.

There is no MAC layer. Scheduling is assumed perfect and collision-free, with no
contention, no hidden terminals, no inter-cluster interference and no capture
effect. Packet loss comes only from distance-dependent
error~\cite{zhao2003understanding,woo2003taming}. Adding interference would
penalize dense-cluster protocols relative to the chain and would make the number
of simultaneously active heads a variable this study cannot see. Latency inherits
the same limitation: it is analytic rather than queueing-based, retransmissions
cost energy but not modeled time, and the absolute figures in
Section~\ref{sec:results} are lower bounds whose ordering is meaningful and whose
magnitudes are not.

Everything is static. Nodes never move and fail only by depletion, so per-link
shadowing is drawn once and held fixed. There is no mobility, no node addition
and no environmental change across the run. The head-to-sink hop is single-hop
for nine of ten configurations, matching each protocol's original specification,
which leaves the chain as the sole multi-hop contrast.

Sensor data is a mean-reverting process rather than a real trace. The reactive
protocols' results are conditional on that model, and a bursty or spatially
correlated field would change them substantially. The thresholds were set
against the process's stationary distribution before any result was observed,
and were not adjusted afterward.

Finally, every protocol targets the same nominal head count but realizes it
differently, and LEACH's realized count is deliberately left stochastic. Part of
every observed difference is therefore attributable to realized cluster count
rather than to selection logic. Fixing the target isolates the question these
protocols actually differ on, namely which five of the roughly five heads
available should be chosen, but it does not eliminate the effect.

\subsection{Scope}

This study establishes results at one energy budget and one traffic model.
Network scale is now swept over three node counts and three field areas, but
within a single family of geometrically similar deployments, and three protocols
have complexity superlinear in node count. Node energy is homogeneous; the
configuration supports heterogeneity but it is disabled throughout. Every node
senses every round, with no duty cycling, event-driven traffic or query
workload. APTEEN's query facilities in particular are out of scope, and only its
reporting-frequency mechanism is modeled. There is no hardware testbed and no
measured trace; every number is simulation output under the assumptions above.

The scale sweep also carries one design limitation worth repeating here. Cells
were chosen to span node count and area independently, not to hold density
constant, so the finding that density does not drive lifetime rests on a close
pairing (50 against 44 nodes per hectare) rather than an exact match. The effect
sizes involved are large enough that a small density gap cannot account for
them, but a grid designed around density would settle the point more directly.

\subsection{On Not Calibrating}

No protocol was tuned to reproduce a published figure. Only protocol
descriptions, the radio model and simulation parameters were taken from prior
work; its numeric results were discarded before implementation began.
Disagreement with published tables is therefore expected and is not evidence of
error here.

The one place this commitment was tested is documented in
Section~\ref{sec:dqn-disclosure}. The DQN's discount factor and training length
were revised after an observed divergence, after which a stopping rule was fixed
in advance and the result recorded as final regardless of outcome. We state this
because the alternative, presenting the final hyperparameters as if they had been
the first, is the ordinary way such revisions disappear from a methods
section~\cite{henderson2018matters,pineau2021reproducibility}.
